\subsection {Cost Allocation to Customers} 
In LRS problems, the aim is to minimize total fixed and variable cost of distributing 
demands to all customers. In real life, besides the total cost, companies also want to 
know the cost of each customer to determine how to charge each customer for the service 
or how to measure profitability of each customer. In literature, people 
investigate for the methods to allocate total cost to each customer considering some
objectives. Objectives might be to be fair for each customer, to find most suitable
cost allocation method for the company, to find a method that is easy to calculate, 
or to be rational. \citet{Tijs} state that not all methods are proper for 
all objectives, thus the company has to decide priority of its objectives before 
choosing a cost allocation method. Cost allocation methods may based on simple
rules or intuition. Also, cooperative game theory is frequently used to define 
cost allocation methods.

We will look at the methods in which game theory is used. In these methods,
a cost game is designed using cost allocation problem. In most of the literature,
the cost game is represented with the pair $(N;c)$, where $N=\{1,2,..,n\}$ is 
the set of players (the customers that we want to allocate cost) and $c$ is the cost 
function defined on $2^N \to \mathbb{R}$. The subsets of N are called coalitions 
and $c(S)$ is the least cost of serving customers in S. Also, it is assumed that 
$c(\varnothing)=0$. $c$ is subadditive if $c(S \cup T) \leq c(S)+c(T)$ for all
$S, T \subset N$ and $S \cap T = \varnothing$. A cost allocation method is a 
function defined for all customers in N and cost function $c$ such that 
$\beta(c)=(\beta_1(c),\beta_2(c),..,\beta_n(c))$ $\in \mathbb{R}^N$ where
$\beta_i(c)$ is the cost assigned to customer i. 

\citet{Tijs} investigate the cost allocation methods based on game theory.  
They have listed properties of the cost allocation methods. For example, if $\beta$ 
function satisfies $\sum_{i \in N}\beta_i(c)=c(N)$, then the method is {\em efficient}.    
The cost allocation method is said to be {\em individually rational} if 
$\beta_i(c) \leq c(\{i\})\ \ \forall i \in N$. The cost allocation method is 
{\em stable} if $\beta(c)$ is in the {\bf core} of the game. The idea of the core is 
an important solution approach for cooperative games. If $x_i$ is the charge for 
player i, then the core of the game is defined as:
\begin{eqnarray}
\label{Core}
CORE(c) = \{ x \in \mathbb{R}^N | \sum_{i \in S}x_i \leq c(S) \ \ \forall S \subset N, \ \ \sum_{i \in N}x_i=c(N)\}               
\end {eqnarray}
$CORE(c)$ is a closed, compact and convex subset of $\mathbb{R}^N$. The core might not 
define a unique point. Also, the core might be empty. Shapley (1971) proved that if c is 
concave, then the core of c is nonempty. c is concave if 
$c(S \cup T)+c(S \cap T) \leq c(S)+c(T)$ $\forall S,T \subseteq N$. if c is concave, the 
game is called convex. Also, it is possible that the core is non-empty when the game 
is not convex.

Finding the core is one of the solution methods for cost allocation problem. Other approach
is finding the {\bf nucleolus}. Nucleolus concept was introduced by \citet{Schmeidler}. 
The nucleolus minimizes maximum dissatisfaction. \citet{Young} 
explains the nucleus concept by comparing two sets S and T. If 
$(c(S)-\sum_{i \in S}x_i) > (c(T)-\sum_{i \in T}x_i)$, S is better than T. To define the
problem that finds the nucleolus, define $e(x,S)=c(S)-\sum_{i \in S}x_i$ that is 
called the excess of S relative to allocation x. To minimize maximum excess, we have to 
solve:
\begin{optprog}
(P1) \qquad \qquad  Max & \objective{w_1} \nonumber \\
subject to & e(x,S) &\geq & w_1 & \forall S \subset N \label{excess-const}\\
           & \sum_{i \in N}x_i &=&c(N)
\end{optprog} 
Note that we have $2^{N-2}$ excess values ($e(x,S)$). By ordering these values from lowest
to highest define vector $e(x)$. \citet{Schmeidler} defined the nucleolus as the vector x 
that maximizes $e(x)$ lexicographically. From the definition of lexicographic order 
(mathematical programming glossary), the vector $e(x)$ is lexicographically greater than 
the vector $e(y)$ if $e(x)-e(y)$ is lexicographically positive that means that the first 
non-zero coordinate of the vector $e(x)-e(y)$ is positive. Every two vectors are either 
equal, or one is lexicographically greater than the other. \citet{Schmeidler} proved that 
there is always a unique vector x (nucleolus) that maximizes e(x). If the core defines a 
non-empty region, then the nucleolus is in the core. \citet{Young} states that the nucleolus 
is the point inside the core that is as far as possible from the boundaries of the core. 
The variations of nucleolus concept are applied in the literature. \cite{Grotte} 
introduced the idea of normalized nucleolus. He defined the excess as 
$e(x,S)=[c(S)-\sum_{i \in S}x_i] / |S|$. \citet{Engevall} defined the demand 
nucleolus with $e(x,S)=[c(S)-\sum_{i \in S}x_i]\sum_{i \in S}D_i$ where $D_i$ is 
the demand of customer i. 

\citet{Kopelowitz} (see also \citet{Dragan}) introduced a method that solves sequence of 
linear programms to find the nucleolus of the game. \citet{Engevall} explains the 
details of the procedure. Procedure starts with the solution of (P1). If (P1) has a 
unique primal solution $(x^*,w_1^*)$, then $x^*$ is the nucleolus of the game. If this is
not the case, by investigating the values of dual variables for constraint 
(\ref{excess-const}), the sets for which the constraint (\ref{excess-const}) is 
binding and active are determined. Let $\pi_1^*(S)$ be the dual variable at optimal 
solution for each $S \subset N$ corresponding to the constraint (\ref{excess-const}). 
If $\pi_1^*(S)>0$, then the constraint (\ref{excess-const}) for set S is binding. By using 
this information, $2^{nd}$ linear programming is defined:
\begin{optprog}
(P2) \qquad \qquad  Max & \objective{w_2} \nonumber \\
subject to & e(x,S) &\geq & w_2 & \forall S \in \{S\subset N | \pi_1^*(S)=0\} \nonumber\\
           & e(x,S) & = & w_1^* & \forall S \in \{S\subset N | \pi_1^*(S)>0\} \nonumber\\
           & \sum_{i \in N}x_i &=&c(N)\nonumber
\end{optprog} 
If P2 has a unique primal solution ($x^*,w_2^*$), then, $x^*$ is the nucleolus, otherwise,
the procedure continues in a similar way by defining P3 and so on. If we define 
$K_q=\{S\subset N| \pi_q^*(S) >0\}$ (the subsets of N such that $w_q^*=e(x,S)$ in 
linear programm at step q), then the problem in
step t is:
\begin{optprog}
($P_t$) \qquad \qquad  Max & \objective{w_t} \nonumber \\
subject to & e(x,S) &\geq & w_t & \forall S \subset N, S \notin \cup_{q=1}^{t-1}K_q \nonumber\\
           & e(x,S) & = & w_q^* & \forall S \in K_q, q=1,..,t-1 \nonumber\\
           & \sum_{i \in N}x_i &=&c(N)\nonumber
\end{optprog} 
The {\bf Shapley value} (introduced by \citet{Shapley}) is another approach for cost 
allocation problems. In this method, players are joining some coalitions in some
random order. A player is assigned the incremental cost that is the cost increase  
for the coalition that the player is attending. The incremental cost for the player changes 
according to the order that he attends. Thus, assuming all orders have same probability
to occur, the expected cost for player i is found to be the following equation.
\begin{eqnarray}
\label{shapley-value}
x_i = \sum_{S \subseteq N | i \in S} \frac{(|S|-1)!(|N|-|S|)!}{|N|!}(c(S)-c(S\backslash \{i\})              
\end {eqnarray}
\citet{Young} interprets Shapley value as the average marginal contribution of each player to
the game if the coalition is obtained by adding one player in every step. The Shapley 
value is a unique solution for the cost allocation, but \citet{Engevall} states that
Shapley value may not be in the core of the game. 

People use these cost allocation methods for practical applications. For example, \citet{Littlechild}
calculate aircraft landing fees and suggest new allocations, \citet{Billera} calculate telephone billing 
rates. In 1998 \citeauthor{Engevall} designed a traveling 
salesman game to find cost allocation of a gas and oil company. They use the minimum cost of  
hamiltonian cycle for the characteristic function (c) of the game. They prove that the core 
of the game might be empty. They calculate the nucleolus and demand nucleolus of the game and 
Shapley value for each customer. Also, they use cost gap allocation method introduced by \citet{Tijs} 
to calculate the cost allocation. They compare the values with the costs
assigned by the company. 

In 2004, \citeauthor{Engevall2} design a vehicle routing game for the 
same gas and oil company in \citet{Engevall}. They use the minimum cost of VRP  
as a characteristic cost function for the game. They use set partitioning 
formulation of VRP. They investigate wheather the core of the game is empty or not. If the core
is not empty, they find a solution in the core.  %G\"othe-Lundgren
\citet{Gothe-Lundgren} prove that
the core of a vehicle routing game is non-empty if and only if LP relaxation of the set 
partitioning formulation of VRP and optimal value of VRP over N customers are equal. Since the 
feasible subsets of N customers are in exponential number, \citet{Engevall2} use a constraint 
generation procedure to determine wheather the core (set \ref{Core}) 
has a solution or not. They use constraint generation procedure to define new set $S \in N$ 
to add the new constraint $\sum_{i \in S}x_i \leq c(S)$ to constraint set of the core.
Additionally, they define an algorithm that uses constraint generation to calculate the 
nucleolus of the game.

In our research, we will define the vehicle routing and scheduling problem as a game similar to  
the approaches in \citet{Engevall,Engevall2}. We will investigate properties
of the game (convex, proper, monotone) and the characteristic function. We will define the 
core of the game and determine wheather the core is empty or not. We will determine the cost
of each customer by using different methods such as simple rules of tumb, Shapley value,   
core and the nucleolus of the vehicle routing and scheduling game. 

